%% ===========================================================================
%%  Interactive Airflow — CHI 2027 submission
%%
%%  TARGET:   ACM CHI 2027, Papers track.
%%            Deadline 10 September 2026, Anywhere on Earth.
%%            Revise & Resubmit 3 December 2026. Conference Pittsburgh, May 2027.
%%
%%  FORMAT:   CHI 2027 requires a SINGLE-COLUMN format for review and an
%%            ANONYMIZED submission, so the class options below are
%%            [manuscript,review,anonymous]. The real author block is kept
%%            below in a comment — uncomment it for the camera-ready and
%%            switch the class options to [sigconf].
%%
%%  LENGTH:   5,000–8,000 words encouraged (excluding references). Over 12,000
%%            is desk-rejected unless justified.
%% ===========================================================================

\documentclass[manuscript,review,anonymous]{acmart}

%% --- CAMERA-READY: swap the line above for the line below -------------------
% \documentclass[sigconf]{acmart}

\usepackage{booktabs}
\usepackage{graphicx}
\usepackage{amsmath}

%% CHI is published under the ACM open-access / rights framework. These are
%% placeholders; the real values arrive from the ACM rights form on acceptance.
\copyrightyear{2027}
\acmYear{2027}
\setcopyright{acmlicensed}
\acmConference[CHI '27]{Proceedings of the 2027 CHI Conference on Human Factors
  in Computing Systems}{May 10--14, 2027}{Pittsburgh, PA, USA}
\acmBooktitle{Proceedings of the 2027 CHI Conference on Human Factors in
  Computing Systems (CHI '27), May 10--14, 2027, Pittsburgh, PA, USA}
\acmDOI{10.1145/XXXXXXX.XXXXXXX}
\acmISBN{978-1-4503-XXXX-X/27/05}

\begin{document}

\title{Saying What You Want to the Air: Turning Non-Expert Comfort Goals into
  Physical Airflow Objectives}

%% ---------------------------------------------------------------------------
%%  AUTHORS.
%%
%%  The real author block lives in authors.tex, which is NOT part of the
%%  submission. CHI has authors upload the LaTeX SOURCE to PCS alongside the
%%  PDF, so names left in a comment here would travel with the submission even
%%  though they never appear in the typeset document. Keeping them in a separate
%%  file means the anonymized build cannot leak them by accident.
%%
%%  SUBMISSION  : leave the \input commented, as below. acmart's `anonymous`
%%                option prints "ANONYMOUS AUTHOR(S)".
%%  CAMERA-READY: switch the class to [sigconf], uncomment the \input, and ship
%%                authors.tex with the source.
%% ---------------------------------------------------------------------------

% %% ===========================================================================
%%  CAMERA-READY AUTHOR BLOCK — DO NOT UPLOAD WITH AN ANONYMOUS SUBMISSION.
%%
%%  CHI 2027 review is anonymous and authors submit the LaTeX source to PCS as
%%  well as the PDF, so this file stays out of the submission bundle entirely.
%%  On conditional acceptance, switch main.tex to \documentclass[sigconf]{acmart},
%%  uncomment its %% ===========================================================================
%%  CAMERA-READY AUTHOR BLOCK — DO NOT UPLOAD WITH AN ANONYMOUS SUBMISSION.
%%
%%  CHI 2027 review is anonymous and authors submit the LaTeX source to PCS as
%%  well as the PDF, so this file stays out of the submission bundle entirely.
%%  On conditional acceptance, switch main.tex to \documentclass[sigconf]{acmart},
%%  uncomment its %% ===========================================================================
%%  CAMERA-READY AUTHOR BLOCK — DO NOT UPLOAD WITH AN ANONYMOUS SUBMISSION.
%%
%%  CHI 2027 review is anonymous and authors submit the LaTeX source to PCS as
%%  well as the PDF, so this file stays out of the submission bundle entirely.
%%  On conditional acceptance, switch main.tex to \documentclass[sigconf]{acmart},
%%  uncomment its \input{authors}, and include this file.
%%
%%  Affiliation confirmed by the author: "Virginia Tech" is the form ACM indexes
%%  and the one used across the Digital Library. The legal name, Virginia
%%  Polytechnic Institute and State University, is also correct but is rarely
%%  used in a CHI author block.
%%
%%  TODO before camera-ready: confirm co-author ordering with the advisors, and
%%  add ORCIDs (\orcid{...}) — ACM asks for them and TAPS will flag their
%%  absence.
%% ===========================================================================

\author{Joseph Leonard}
\affiliation{%
  \institution{Virginia Tech}
  \city{Blacksburg}
  \state{Virginia}
  \country{USA}}
\email{josephl04@vt.edu}

\author{Haoran Xie}
\affiliation{%
  \institution{Japan Advanced Institute of Science and Technology}
  \city{Nomi}
  \state{Ishikawa}
  \country{Japan}}

\author{Takeo Igarashi}
\affiliation{%
  \institution{The University of Tokyo}
  \city{Tokyo}
  \country{Japan}}

\author{Bo Zhu}
\affiliation{%
  \institution{Georgia Institute of Technology}
  \city{Atlanta}
  \state{Georgia}
  \country{USA}}
, and include this file.
%%
%%  Affiliation confirmed by the author: "Virginia Tech" is the form ACM indexes
%%  and the one used across the Digital Library. The legal name, Virginia
%%  Polytechnic Institute and State University, is also correct but is rarely
%%  used in a CHI author block.
%%
%%  TODO before camera-ready: confirm co-author ordering with the advisors, and
%%  add ORCIDs (\orcid{...}) — ACM asks for them and TAPS will flag their
%%  absence.
%% ===========================================================================

\author{Joseph Leonard}
\affiliation{%
  \institution{Virginia Tech}
  \city{Blacksburg}
  \state{Virginia}
  \country{USA}}
\email{josephl04@vt.edu}

\author{Haoran Xie}
\affiliation{%
  \institution{Japan Advanced Institute of Science and Technology}
  \city{Nomi}
  \state{Ishikawa}
  \country{Japan}}

\author{Takeo Igarashi}
\affiliation{%
  \institution{The University of Tokyo}
  \city{Tokyo}
  \country{Japan}}

\author{Bo Zhu}
\affiliation{%
  \institution{Georgia Institute of Technology}
  \city{Atlanta}
  \state{Georgia}
  \country{USA}}
, and include this file.
%%
%%  Affiliation confirmed by the author: "Virginia Tech" is the form ACM indexes
%%  and the one used across the Digital Library. The legal name, Virginia
%%  Polytechnic Institute and State University, is also correct but is rarely
%%  used in a CHI author block.
%%
%%  TODO before camera-ready: confirm co-author ordering with the advisors, and
%%  add ORCIDs (\orcid{...}) — ACM asks for them and TAPS will flag their
%%  absence.
%% ===========================================================================

\author{Joseph Leonard}
\affiliation{%
  \institution{Virginia Tech}
  \city{Blacksburg}
  \state{Virginia}
  \country{USA}}
\email{josephl04@vt.edu}

\author{Haoran Xie}
\affiliation{%
  \institution{Japan Advanced Institute of Science and Technology}
  \city{Nomi}
  \state{Ishikawa}
  \country{Japan}}

\author{Takeo Igarashi}
\affiliation{%
  \institution{The University of Tokyo}
  \city{Tokyo}
  \country{Japan}}

\author{Bo Zhu}
\affiliation{%
  \institution{Georgia Institute of Technology}
  \city{Atlanta}
  \state{Georgia}
  \country{USA}}


\renewcommand{\shortauthors}{Anonymous et al.}

%% ===========================================================================
\begin{abstract}
People describe indoor air in the language of comfort: a bedroom that is too
cold, a kitchen smell that reaches the bed, a draft that wakes them. Tools that
could answer those complaints — real-time fluid solvers and inverse-design
methods — instead require the problem stated in physics: a target temperature in
a named zone, a supply velocity at a diffuser, an air-change rate. Prior systems
assume that statement already exists. This paper asks whether it can be
\emph{produced} from what a non-expert would actually say.

We present \textsc{Interactive Airflow}, a browser-based 3D home editor in which
a comfort goal expressed by hand, in typed language, or as a drawing is grounded
onto a small closed vocabulary of physical objectives, evaluated live by an
in-browser Euler solver, and handed to a constrained optimizer that returns
complete furniture-and-device arrangements. The closed vocabulary is the
mechanism that keeps the translation checkable: every interpreted goal is a
tuple the solver can score, so a suggestion can never be plausible-sounding and
physically wrong. The optimizer treats placement as a constrained combinatorial
problem over a discretized design space, solved with a three-rung multi-fidelity
cascade and coordinate-descent refinement; against exhaustive search at
reporting fidelity it lands in the top 0.8\% and top 1.4\% of the feasible space
on two tasks while returning in 3--8\,s in the browser.

We report a within-subjects study with 15 non-experts across four residential
scenarios, comparing manual editing against the full tool. \textbf{[Results
pending analysis --- see Section~\ref{sec:results}.]}
\end{abstract}

%% CCS concepts. Generate the official block at
%% https://dl.acm.org/ccs and paste it here before submission.
\begin{CCSXML}
<ccs2012>
<concept>
<concept_id>10003120.10003121.10003122</concept_id>
<concept_desc>Human-centered computing~HCI design and evaluation methods</concept_desc>
<concept_significance>500</concept_significance>
</concept>
<concept>
<concept_id>10003120.10003121.10003124.10010870</concept_id>
<concept_desc>Human-centered computing~Natural language interfaces</concept_desc>
<concept_significance>500</concept_significance>
</concept>
<concept>
<concept_id>10010147.10010371.10010352</concept_id>
<concept_desc>Computing methodologies~Animation</concept_desc>
<concept_significance>300</concept_significance>
</concept>
</ccs2012>
\end{CCSXML}

\ccsdesc[500]{Human-centered computing~HCI design and evaluation methods}
\ccsdesc[500]{Human-centered computing~Natural language interfaces}
\ccsdesc[300]{Computing methodologies~Physical simulation}

\keywords{indoor airflow, natural language interfaces, intent, physical
  simulation, design optimization, non-experts, sketch interaction}

\maketitle

%% ===========================================================================
\section{Introduction}

A person who is uncomfortable at home knows exactly what is wrong and has no
words for the fix. The bedroom is cold in winter while the living room is fine.
The kitchen bin smells and the bed is four metres away. The air conditioner is
aimed down the length of the room at the pillow, and it is the only thing
keeping the flat cool. These are ordinary complaints, stated the way people
state them.

The tools that could resolve them do not accept that statement. Computational
fluid dynamics for indoor environments is mature, and so are the methods built
on top of it: inverse design that adjusts HVAC parameters until a numerical
target is met~\cite{liu2017inverse}, real-time solvers fast enough to sit inside
an interactive loop~\cite{zhu2011sketchfluid}, and learned surrogates that
replace the solve
entirely~\cite{thuerey2020deep}. Every one of them begins after the hard part is
already done. Inverse design needs \emph{``hold 24\,\textdegree C in zone A,
supply velocity 0.5\,m/s at diffuser D2, draft below 0.2\,m/s.''} Sketch-based
systems need the user to know \emph{where} the sources and sinks belong. Even
recent work applying large language models to air-distribution
design~\cite{dai2025llm} takes an expert-authored prompt that already presumes
CFD knowledge, and asks a professional to correct the output.

The gap is at the input. Between \emph{``keep my bed cool at night, but I hate a
draft on my face''} and the objective vocabulary those systems consume, there is
a translation nobody has built, and it is not a lookup. One comfort goal maps to
many physical configurations, so the system must select objectives rather than
recover a unique answer. Goals conflict in ways the person stating them does not
realize --- air speed cools and is also what a draft is made of --- so the
translation must produce competing objectives, not a single target. Nouns like
``the bed'' and ``the kitchen'' are references that must resolve to regions in
\emph{this} home. And the result has to be verifiable: a language model can emit
fluent, well-formatted, physically wrong numbers, and a system that cannot tell
the difference is worse than no system.

This paper takes that translation as the research problem. Our contributions are:

\begin{enumerate}
  \item \textbf{An intent-to-physics layer} that grounds everyday comfort goals
    onto a closed vocabulary of evaluable objectives --- a tuple of region,
    scalar, direction and hardness --- rather than onto free-form solver
    parameters. Because the vocabulary is closed, every interpretation is
    checkable against the solver that draws the user's own view, which is what
    makes the translation trustworthy rather than merely fluent.
  \item \textbf{Three interchangeable ways to state the same intent} --- direct
    manipulation, typed language, and drawing --- all reducing to the same
    objective tuples, so a study can observe which modality people reach for and
    when without confounding modality with capability.
  \item \textbf{A constrained optimizer} over device placement, aim and openings
    that returns complete arrangements a participant could have reached by hand,
    with feasibility enforced by construction and task rules treated as hard
    constraints. We characterize it against exhaustive search.
  \item \textbf{An empirical study} with 15 non-experts over four residential
    scenarios, contrasting manual editing with the full tool, and an instrumented
    session log built for behavioral analysis rather than for reading.
\end{enumerate}

The system runs entirely in a browser with no installation, no server and no
account, which is what made a remote study of non-experts practical at all.

%% ===========================================================================
\section{Related Work}

\subsection{Inverse design of indoor environments}

Liu et al.~\cite{liu2017inverse} develop an adjoint method over a fast
fluid-dynamics solver for the inverse design of indoor environments: the
designer states a numerical target, the method computes how HVAC parameters
influence it, and iterates. The machinery is exactly what a comfort complaint
would need --- and it starts from a physical objective an occupant cannot write.
Our work is deliberately the step \emph{before} this one: producing the
objective that inverse design consumes.

\subsection{Sketch-based fluid authoring}

SketchFluid~\cite{zhu2011sketchfluid} lets users draw a fluid system --- pipes,
regions, sources, sinks --- on a 2.5D canvas, with a hydraulic and 2D
Navier--Stokes solver filling in plausible flow interactively, demonstrated on
indoor air with multiple air conditioners. Drawing is a far lower barrier than
numerical targets, but it still asks the user to know where the sources and
sinks go: the answer is the input. This work is a direct descendant --- sketch
to language, 2.5D to 3D --- and retains drawing as one modality among three
precisely because sketch and language fail in different places.

\subsection{Learned surrogates for flow}

Thuerey et al.~\cite{thuerey2020deep} train a U-Net to predict RANS solutions
for airfoil flows in a single forward pass at under 3\% mean error. Surrogates
make simulation fast; they do not make it addressable. Speed is enabling
technology for the interactive loop we need, not a contribution of this work.

\subsection{Language models for HVAC design}

Dai et al.~\cite{dai2025llm} apply GPT-4o to office air-distribution design: an
HVAC professional writes a structured prompt, the model proposes a design, the
professional corrects it, and one CFD run checks the result. This establishes
that using a language model in this domain is not itself novel. The distinction
our work draws is between \emph{accepting natural language} and \emph{accepting
a non-expert's comfort goal}: their prompt is expert-authored and already
physical. Table~\ref{tab:related} makes the comparison explicit; the column no
prior system fills is the last one.

\begin{table}[t]
  \caption{Prior work by capability. Every system assumes the user can already
    supply a physical objective, a correct sketch, or an expert prompt. The
    empty column is the contribution.}
  \label{tab:related}
  \small
  \begin{tabular}{lccccc}
    \toprule
    & Indoor & Real-time & Optimi- & NL & Non-expert \\
    & airflow & sim & zation & input & comfort goals \\
    \midrule
    Dai et al.~\cite{dai2025llm}       & \checkmark & --- & --- & \checkmark & --- \\
    SketchFluid~\cite{zhu2011sketchfluid} & \checkmark & \checkmark & --- & --- & --- \\
    Deep-Flow~\cite{thuerey2020deep}   & --- & \checkmark & --- & --- & --- \\
    Inverse design~\cite{liu2017inverse} & \checkmark & $\sim$ & \checkmark & --- & --- \\
    \textbf{This work}                 & \checkmark & \checkmark & \checkmark & \checkmark & \checkmark \\
    \bottomrule
  \end{tabular}
\end{table}

%% ===========================================================================
\section{Formative Study and Design Rationale}
\label{sec:formative}

Before building the interaction we ran a formative questionnaire with
non-experts ($n=8$), following a story-first, opinions-last
structure~\cite{mackay2023doit}, to learn what people would actually do about an
uncomfortable room.

\subsection{People reach for a device; openings are an afterthought}

Asked to rank what they would do first to cool a room quickly, respondents
reached for hardware and treated the building itself as inert
(Table~\ref{tab:formative}). Not one would open a window first.

\begin{table}[t]
  \caption{``To cool a room quickly, rank what you would do first'' ($n=8$).
    Lower is earlier.}
  \label{tab:formative}
  \small
  \begin{tabular}{lcc}
    \toprule
    Action & Mean rank & Ranked first by \\
    \midrule
    Turn on the AC   & 1.9 & 5 / 8 \\
    Turn on a fan    & 1.9 & 2 / 8 \\
    Open the door    & 3.1 & 1 / 8 \\
    Open windows     & 3.1 & \textbf{0 / 8} \\
    \bottomrule
  \end{tabular}
\end{table}

The solver disagrees. In our example home at 33\,\textdegree C outdoors with the
air conditioner running in the living room, opening one interior door moves the
bedroom from 33.0\,\textdegree C to 25.8\,\textdegree C while the living room is
unchanged at 22.9 versus 22.8\,\textdegree C --- 7.2\,\textdegree C for free.
Winter is symmetric: with a heater in the living room at 2\,\textdegree C
outdoors, the bedroom goes from 2.0 to 7.4\,\textdegree C when the door opens.

This gap between what people would do and what works is the reason the system
exists, and it shaped every task in the study: each is constructed so that the
device-only answer is insufficient and an opening or an aim is the real lever.
It also yields a falsifiable form of ``is the tool useful'':

\begin{quote}
  Does the tool move non-experts off the device-first strategy and onto the
  configuration the physics supports --- and do they believe it when it does?
\end{quote}

\subsection{Design goals}

\begin{description}
  \item[D1. Never require a physical quantity.] Every input the participant
    gives is a comfort goal, a direct manipulation, or a drawing. Thresholds
    exist and are never shown: a participant told ``keep it under
    0.20\,m/s'' optimizes toward a number instead of judging whether the room
    feels right, and the number is the thing the tool is supposed to be
    explaining.
  \item[D2. Every interpretation must be checkable.] Intent is grounded onto a
    closed objective vocabulary that the solver can score, so an interpretation
    can be right or wrong rather than merely plausible.
  \item[D3. Suggestions must be reachable by hand.] The optimizer may only
    propose changes the participant is permitted to make. A card that relocates
    a bolted vent teaches the wrong lesson and cannot be undone by the person
    holding the mouse.
  \item[D4. One solve behind everything.] The visualization, the goal checks and
    the optimizer's final ranking read the same solve, so what the participant
    sees and what the system claims cannot disagree.
\end{description}

%% ===========================================================================
\section{\textsc{Interactive Airflow}}

%% ---------------------------------------------------------------------------
%%  FIGURES — none are placed yet. Three are needed before submission:
%%
%%   1. TEASER (page-width, top of page 1): the three input modalities side by
%%      side on the same home — a drag, a typed sentence, a drawn box+arrow —
%%      each with the layout it produced.
%%   2. SYSTEM PIPELINE (column-width): intent -> objective tuple -> solver ->
%%      optimizer -> gallery.
%%   3. BEFORE/AFTER FIELD (page-width): the apartment task with the louvre
%%      level vs aimed down, streamlines coloured by temperature.
%%
%%  CHI RULES THAT APPLY TO ALL OF THEM:
%%   - \Description{} is REQUIRED on every figure (SIGCHI accessibility guide).
%%     It is not the caption: it must carry the information a sighted reader
%%     gets from the image and that the caption does not already state. It does
%%     not appear in the PDF; TAPS puts it in the HTML5 version.
%%   - Size figures at their FINAL 2-column size. A column-wide figure should
%%     appear half-page-width in this 1-column submission, leaving white space
%%     to its right. Do not scale images — export at the exact size.
%%   - Do not encode meaning in colour alone; add shape or pattern.
%%
%%  Template:
%%
%%  \begin{figure}[t]
%%    \centering
%%    \includegraphics[width=0.5\textwidth]{figs/teaser}
%%    \caption{Short, states what the reader should take from it.}
%%    \Description{Full description for a reader who cannot see the image.}
%%    \label{fig:teaser}
%%  \end{figure}
%% ---------------------------------------------------------------------------

\subsection{Walkthrough}

The participant opens a prebuilt home in a 3D editor and is given a comfort goal
in plain words. They may drag devices, aim them, open and close doors and
windows, and run the simulation as often as they like. Flow is drawn as
temperature-colored streamlines over a floor-level field for temperature,
humidity or contaminant. When they are satisfied they press Submit, which scores
the goals a final time and downloads the session as JSON.

\subsection{Three ways to state the same intent}

\textbf{By hand.} Direct manipulation of devices and openings on a 0.25\,m grid,
with collision and doorway keep-clear rules enforced during the drag.

\textbf{By typing.} A curated domain dictionary resolves the common vocabulary
of comfort offline and instantly --- \emph{cool} to minimizing temperature,
\emph{stuffy} to increasing ventilation, \emph{smell} to minimizing contaminant
transport from a source region, \emph{drafty} to capping air speed. A language
model is consulted only for wording the dictionary cannot match, and is
constrained to emit \emph{only} tuples in the same vocabulary, which is what
keeps its output verifiable. A sentence neither can read still produces a
search, run against the active task's own goals, and says so rather than failing
silently.

\textbf{By drawing.} The participant boxes a region and marks it \emph{warm},
\emph{cool}, \emph{fresh air} or \emph{no wind}, or draws an arrow to move air
from one place to another. Regions resolve against the editor's scene graph.

All three reduce to the same objective tuples, and downstream nothing knows
which was used.

\subsection{The objective vocabulary}

Every goal becomes a tuple
\[
  \textit{objective} = (\textit{region},\ \textit{scalar},\ \textit{direction},\
  \textit{hard} \mid \textit{soft})
\]
with \textit{scalar} drawn from a closed set --- temperature, air speed,
ventilation, contaminant concentration, drying time --- and \textit{direction}
one of minimize, maximize, or hold near target. Hard tuples are constraints that
must hold; soft tuples are objectives to improve. The vocabulary is deliberately
small. It is what allows a typed sentence, a drawn box, and a task's own
built-in goals to be scored by exactly the same machinery, and what lets the
system report that it did not understand rather than inventing a target.

\subsection{Simulation}
\label{sec:sim}

Flow is computed by an incompressible Euler solver written in TypeScript and
running entirely in the browser. The scheme is the standard \emph{Stable Fluids}
formulation~\cite{stam1999stable}: unconditionally stable semi-Lagrangian
advection followed by a pressure projection~\cite{chorin1968projection}, on a
staggered marker-and-cell grid~\cite{harlow1965mac} with velocity components on
cell faces and pressure and scalars at cell centres. We adopt it for the reason
it was introduced --- stability at large time steps, independent of the CFL
condition --- because the interaction loop needs a steady field in well under a
second and cannot afford sub-stepping. Buoyancy is added as a temperature-driven
body force in the manner of Fedkiw et al.~\cite{fedkiw2001visual}, which is what
makes the third dimension necessary: a top-down slice cannot represent warm air
rising or cold air pooling, and both are the whole subject of the tasks.

\textbf{Discretization of the home.} The plan is voxelized into cells labelled
solid (walls, furniture), open (an exterior opening, treated as a free boundary
through which air may leave the domain), or fluid. At reporting fidelity the
four study homes resolve to $18{\times}8{\times}21$ to $19{\times}14{\times}16$
cells --- 3{,}024 to 4{,}256 in total, at a cell size of 0.24--0.40\,m. This is
a coarse grid by any CFD standard and is a deliberate consequence of the
interaction budget, not an oversight; Section~\ref{sec:technical} reports what
it is and is not calibrated for.

\textbf{Devices as boundary conditions.} A step advances as: apply body forces
and buoyancy, advect velocity semi-Lagrangian by backtracing each face and
interpolating trilinearly, project to a divergence-free field by Gauss--Seidel
relaxation with over-relaxation, then advect the scalars. Devices enter this
loop in two physically distinct ways, and the distinction turned out to matter
for how participants reason about them. A vent or an air conditioner
\emph{prescribes} a velocity on its outgoing faces: it supplies air at a set
speed regardless of what the room does. A fan instead applies a \emph{body
force} to the air already in front of it, and is therefore subject to the same
pressure projection as everything else --- so it weakens when it pushes into a
flow that resists, and it cannot manufacture air where there is none. A fan
aimed at a closed door does very little, which is the correct answer and the one
a prescribed-velocity model would get wrong.

\textbf{Transported scalars.} Temperature, contaminant concentration and drying
are carried along the converged velocity field by a geodesic transport pass,
rather than being advected to convergence themselves. Exterior glazing chills
the air against it, so the cold downdraft that pools under a window in winter
forms for the reason it does in a real room.

\textbf{Asymmetry of a cooling jet.} One modelling detail deserves its own
statement because the apartment task turns on it entirely. Supply air from an
air conditioner is roughly 8--10\,K below room temperature and therefore denser
than the air around it, so the louvre is either working with gravity or against
it. Aimed down, the jet keeps its momentum to the floor, lands, and spreads
outward as a gravity current that reaches corners direct throw never would.
Aimed up, it decelerates against its own weight within a metre, stalls, and
falls back over the unit that threw it --- the short-circuiting that is the
reason no cooling louvre is aimed at a ceiling. An earlier version of the model
scaled the jet symmetrically in tilt, which made the two indistinguishable and
gave the task two contradictory answers; the current model penalizes only the
upward case, and the difference between a louvre that cools the far room and one
that does not is 24.7\,\textdegree C against 25.6\,\textdegree C in the warmest
corner.

\subsection{The optimizer}
\label{sec:optimizer}

Intent becomes a constrained combinatorial optimization over device placement,
solved against the same solver that draws the views.

\textbf{Decision variables.} Per movable device: position on a 0.25\,m lattice,
heading, and vertical aim quantized to 15\textdegree{} over $\pm$60\textdegree{}
($\pm$30\textdegree{} for free-standing units). Per opening, a boolean. A
wall-mounted unit contributes aim only, its heading being fixed by the wall it is
bolted to. Feasibility is enforced by construction rather than by rejection:
candidate positions are projected onto free space by a collision solver that
already knows about furniture, doorway keep-clear zones and per-task room
confinement, so infeasible points are never scored.

\textbf{Objective.} A scalarization of the parsed request over the solver's
steady fields, with the active task's thresholds entering as constraints. The
composite is lexicographic in three tiers:
\[
  \textit{score} = \textit{request}(x) - 100\,|\mathcal{V}(x)| - 3\sum_i s_i(x) -
  \sum_i m_i(x)
\]
where $\mathcal{V}$ is the set of violated constraints, $s_i$ the shortfall on
goal $i$ normalized by that goal's own tolerance --- so a 0.03\,m/s draft
overshoot and a 0.4\,\textdegree C temperature overshoot are commensurable ---
and $m_i$ the same distance measured against a bar drawn a quarter-tolerance
inside the real one. The last term orders solutions that all satisfy the
constraints by how much headroom they leave, so nothing is recommended sitting
exactly on its own limit. The tiering matters: with a purely linear penalty the
optimizer will sell a violated constraint for a good number elsewhere, which on
the apartment task meant trading 0.22\,m/s across a sleeper's bed for a degree in
the living room.

\textbf{Multi-fidelity cascade.} Evaluating every candidate at reporting fidelity
is far too slow for an interactive loop, so solver cost is spent where it
discriminates (Table~\ref{tab:cascade}). Each rung requests a target cell count
and a number of relaxation iterations and time steps; the target is realized as
whichever grid best fits the home, so the delivered resolution varies by task.
Because the final rung is the fidelity the goal checks themselves use, a
suggestion's printed numbers and its checkmarks cannot disagree (D4).

\begin{table}[t]
  \caption{The three rungs of the cascade, with the grids they resolve to across
    the four study homes.}
  \label{tab:cascade}
  \small
  \begin{tabular}{lrrrr}
    \toprule
    Rung & Target & Iters & Steps & Delivered cells \\
    \midrule
    Screen & 1{,}200 & 4 & 8  & 960--1{,}260 \\
    Mid    & 1{,}800 & 6 & 12 & 1{,}155--1{,}890 \\
    Final  & 4{,}200 & 8 & 22 & 3{,}024--4{,}256 \\
    \bottomrule
  \end{tabular}
\end{table}

The middle rung is not always a distinct rung. On two of the four homes the
screen and mid targets round to the same grid, so the re-rank there differs from
the screen only in iteration and step count; the cascade is genuinely three-level
on the other two. We report this rather than describing a uniform ladder,
because it bounds how much of the accuracy in Section~\ref{sec:technical} can be
attributed to the middle stage.

\textbf{Local refinement.} Finalists are polished by coordinate descent over
$(x, z, \textit{heading}, \textit{tilt})$ --- eight translations, two rotations
and two tilts per step, each projected back onto the feasible set --- under a
separate 70-evaluation budget. This is what reaches positions off the candidate
lattice.

\textbf{Constraints from the drawing.} A sketched arrow is treated as a
directional constraint rather than an objective term: it restricts admissible
headings to a 60\textdegree{} cone about the bearing to the arrow's head, so the
returned layout satisfies the drawn direction by construction instead of trading
it away against the scalar objective.

%% ===========================================================================
\section{Technical Evaluation}
\label{sec:technical}

\subsection{Optimizer quality against exhaustive search}

We brute-forced two tasks over their reachable space at reporting fidelity and
compared the result with what the optimizer returns
(Table~\ref{tab:optimizer}).

\begin{table}[t]
  \caption{Optimizer against exhaustive search at reporting fidelity.}
  \label{tab:optimizer}
  \small
  \begin{tabular}{llrrr}
    \toprule
    Task & Request & Start & Exhaustive & Optimizer \\
    \midrule
    Winter & \emph{warm the bedroom} & 13.51 & 21.49 & \textbf{20.77} \\
    Apartment & \emph{cool the flat} & 24.89 & 24.58 & \textbf{24.66} \\
    \bottomrule
  \end{tabular}
  \\[2pt]
  \footnotesize All values \textdegree C. Winter: top 0.8\% of 384 layouts.
  Apartment: top 1.4\% of the 148 \emph{feasible} layouts.
\end{table}

The apartment row is scored against the feasible subset deliberately. The
unconstrained optimum over that grid is 24.47\,\textdegree C and it drives
1.15\,m/s across the bed --- six times the task's draft limit. It is the coldest
arrangement in the space and it is not an answer to the task; 284 of the grid's
432 layouts are infeasible in the same way. Since the optimizer solves the
constrained problem, the honest comparison is against the 148 layouts that also
keep the bed calm, and there it ranks third. A search completes in roughly
3--8\,s single-threaded in the browser.

\subsection{Feasibility auditing}

Design goal D3 is verified rather than assumed. On every change to the search,
all four tasks are run and every returned layout is diffed against what its task
permits --- no relocating a bolted vent, no turning a wall-mounted unit, no
touching a power dial the task fixes. The audit is what caught, among others, a
suggestion that switched off a fan on a task whose power control is hidden,
leaving the participant unable to switch it back on.

\subsection{Threshold calibration}

Every task threshold was set by measurement rather than by judgement, by
bundling the simulation modules headlessly and sweeping the reachable space. The
humidity task's 66-minute drying line, for example, separates 57 of 104 vent and
window combinations from the rest, and every arrangement leaving the extract on
the window's own wall reads 113 minutes or worse whatever the glazing does. The
sweep that produced each number is recorded beside it in the source.

%% ===========================================================================
\section{User Study}
\label{sec:study}

\subsection{Design}

A within-subjects study contrasting \textbf{manual editing} with the
\textbf{full tool} (typed and drawn intent plus the optimizer), with condition
order counterbalanced across participants and matched scenario sets assigned so
that neither condition is consistently paired with the same scenarios.

\subsection{Participants}

Fifteen non-experts (8 male, 7 female), spanning a small number of teenagers, a
majority in their twenties, and a small number in their thirties. Anyone with
HVAC, CFD, mechanical-engineering or architecture training was screened out, the
claim being about non-experts. Sessions ran remotely over video with
screen-sharing, approximately 45--50 minutes each.

\textbf{[TODO(Joseph): confirm the age bands as numbers, the recruitment
channel, the compensation, and the ethics approval reference for the ethics
statement below.]}

\subsection{Tasks}

Four residential scenarios, each a prebuilt home, a fixed set of permitted
controls, and one or two checkable goals held in tension so that no single move
finishes the job (Table~\ref{tab:tasks}). Task wording gives the comfort goal
only --- never a physical quantity, never a device, never a room-to-room
mechanism (D1).

\begin{table}[t]
  \caption{The four study tasks. Every goal is a solver-checked threshold that
    the participant is never shown.}
  \label{tab:tasks}
  \small
  \begin{tabular}{p{1.7cm}p{2.2cm}p{3.4cm}}
    \toprule
    Task & May change & Goal \\
    \midrule
    Winter & heater (living room only), fan & both rooms comfortable, no cold pool at the glass \\
    Kitchen smell & fan, either window & keep the kitchen smell off the bed \\
    Humidity & extract vent, window & dry the bathroom fast after a shower \\
    AC on the bed & louvre tilt, fan, bedroom door & no strong draft on the bed \emph{and} both rooms cool everywhere \\
    \bottomrule
  \end{tabular}
\end{table}

\subsection{Measures}

\textbf{Goal attainment.} Solver-scored constraint satisfaction on the final
state --- binary per objective, plus the achieved margin.

\textbf{Strategy shift.} Whether the final state changes at least one opening,
and edit counts by object type, testing directly whether participants move off
the device-first strategy of Section~\ref{sec:formative}.

\textbf{Calibration and silent error.} A belief prompt (yes/no plus a five-point
confidence rating) asked \emph{before} any verdict is shown, crossed with the
solver's score --- the cases where a participant believes the goal held and it
did not.

\textbf{Effort.} Time to self-declared completion, number of edits, number of
simulation runs.

\textbf{Modality use.} Which of the three input routes each edit came through,
and in what order across a session.

\textbf{Subjective.} Perceived usefulness for a real room rather than for the
task, and SUS, so that usefulness is not confounded with interface friction.

\subsection{Instrumentation}

Submit seals a JSON record built for analysis rather than for reading. Every
event carries the method it came through (\texttt{manual}, \texttt{text},
\texttt{sketch}, \texttt{solution} or \texttt{system}) and the \emph{whole
layout after the event}, so any step can be scored, diffed or re-simulated on
its own without replaying the ones before it. Simulator readings are recorded
whenever the value moves rather than only when a checkbox flips. Typed sentences
are stored verbatim together with how they were parsed, before anything is done
with them; a search records what it was asked, what it was permitted to touch,
and every card it offered with that card's full layout and predicted numbers, so
an accepted suggestion can be compared against the ones passed over.

\subsection{Analysis plan}

At $n=15$ with one scenario per condition per participant this is a
mixed-methods formative-summative study, not a powered confirmatory trial. We
report medians, per-participant paired differences, effect sizes and confidence
intervals, with Wilcoxon signed-rank tests where a paired comparison is
meaningful, and treat every between-modality comparison as descriptive and
qualitative.

%% ===========================================================================
\section{Results}
\label{sec:results}

\textbf{[PENDING --- the 15 session logs have not yet been analysed. Structure
below; numbers to follow.]}

\subsection{Goal attainment}
\textbf{[Manual vs.\ full tool, per task; paired differences; margins.]}

\subsection{Strategy shift}
\textbf{[Proportion changing an opening, by condition. This is the section that
answers the formative finding in Section~\ref{sec:formative}.]}

\subsection{Calibration}
\textbf{[Belief $\times$ solver score; silent-error rate by condition.]}

\subsection{Modality use}
\textbf{[Counts by method from the session logs; ordering within sessions;
switching behaviour.]}

\subsection{Effort and subjective response}
\textbf{[Time, edits, simulation runs; usefulness and SUS.]}

%% ===========================================================================
\section{Discussion}

\textbf{[PENDING --- to be written against the results.]}

%% ===========================================================================
\section{Limitations}

The solver is a coarse real-time prototype: 3{,}024--4{,}256 cells at reporting
fidelity, a cell size of 0.24--0.40\,m, and an inviscid Euler formulation with
no turbulence model. At that resolution a bed, a doorway and a louvre are a
handful of cells each. The model is calibrated so that the four study tasks
order layouts the way the physics does, and it has not been validated against
physical measurement or against a reference CFD solve; the claims here are about
a consistent interactive model that participants can reason with, not about
predictive accuracy for a real building. Establishing the latter would require a
convergence study against a fine-grid RANS or LES reference, which we have not
run.

The optimizer is a budgeted local method over a discretized design space and
carries no optimality certificate. Its guarantee is empirical, measured against
exhaustive search on two tasks, and the cheap rung of the cascade can in
principle discard a region the accurate rung would have preferred.

Scalar fields beyond temperature, contaminant and drying --- CO\textsubscript{2}
and particulates in particular --- are not implemented, which limits the range of
comfort goals the vocabulary can currently ground.

The study is small and remote, with four scenarios in a single prebuilt home
each. It can show whether the tool changes strategy and belief in these
scenarios; it cannot establish that the effect generalizes to a participant's
own home, which is the setting the work is ultimately about.

The visualization encodes temperature as a diverging colour ramp, which is
information carried by colour alone. The airflow view offers two encodings ---
moving particles or streamlines --- so motion and shape are available there, but
a reader with a colour-vision deficiency has no non-colour route to the
temperature field. A redundant encoding (contour banding, or hatching at the
comfort boundaries) is the obvious fix and is not yet implemented.

\section{Conclusion}

\textbf{[PENDING --- to be written against the results.]}

%% ===========================================================================
\begin{acks}
\textbf{[Anonymized for review. Add funding, the lab, and the pilot participants
for the camera-ready.]}
\end{acks}

\bibliographystyle{ACM-Reference-Format}
\bibliography{refs}

\end{document}
